% concatenated from main.tex
\documentclass[11pt]{article}

\usepackage{amssymb}
\usepackage{amsmath}
\usepackage{amsthm}
\usepackage{mathrsfs}
\usepackage{graphicx}
\usepackage[utf8]{inputenc}
\usepackage{multirow}
\usepackage{comment}
\usepackage{tikz}
\usetikzlibrary{tikzmark}

\setlength{\oddsidemargin}{0.0in}
\setlength{\evensidemargin}{0.0in}
\setlength{\topmargin}{-0.5in}
\setlength{\textheight}{9.5in}
\setlength{\textwidth}{6.5in}

\font\smallit=cmti10
\font\smalltt=cmtt10
\font\smallrm=cmr9

\newtheorem{thm}{Theorem}[section]
\newtheorem{prop}{Proposition}
\newtheorem{lem}{Lemma}[section]
\newtheorem{cor}{Corollary}
\newtheorem{rem}{Remark}
\newtheorem{ex}{Example}
%\newtheorem{def}{Definition}
\newenvironment{Pf}{\noindent{\bf Proof. }}{\hfill $\blacksquare$ \\}

\def\mod{\!\pmod}
\def\g{\tt g}
\def\m{\tt m}
\def\n{\tt n}
\def\s{\tt s}
\def\C{\mathscr C}
\def\G{\Gamma}
\def\S{\mathcal S}
\def\Imp{\Rightarrow}
\def\ds{\displaystyle}
\def\ts{\textstyle}
\def\lf{\lceil}
\def\rf{\rceil}
\def\lc{\lceil}
\def\rc{\rceil}
\def\Iff{\Leftrightarrow}
\def\LIff{\Longleftrightarrow}
\def\D{\Delta}
\def\G{\Gamma}
\def\L{\Lambda}
\def\N{\mathbb N}
\def\Z{\mathbb Z}
\def\a{\alpha}
\def\b{\beta}
\def\k{\kappa}
\def\l{\lambda}
\def\lang{\langle}
\def\rang{\rangle}
\def\d{\delta}
\def\e{\epsilon}
\def\m{\bf m}
\def\s{\bf s}
\def\ov{\overline}
\def\es{\emptyset}
\def\imp{\rightarrow}
\def\pr{\prime}
\def\sm{\setminus}


%%%%%%%%%%%%%%%%%%%%%%%%%%%%%%%%%%%%%%%%%%%%%%%%%%%%%%%
% % Lines starting with % are comments, which are ignored.
% % This is a handy way of indicating the date and version of
% % your document, to wit:
% %
% % LaTeX sample file
% % Modified March, 2002
% %
%%%%%%%%%%%%%%%%%%%%%%%%%%%%%%%%%%%%%%%%%%%%%%%%%%%%%%%
% % Title and author(s)
%%%%%%%%%%%%%%%%%%%%%%%%%%%%%%%%%%%%%%%%%%%%%%%%%%%%%%%
\title{Improved Bounds on  Diffsequences with Gaps in Powers of 2}
\author{Kanav Talwar\thanks{
                  Delhi Public School Faridabad, Faridabad, HR 121002, India}
        \and
        Utkarsh Gupta\thanks{
                      Department of Electrical and Computer Engineering, Northeastern University, Boston, MA 02115, USA
                  }
        }
%%%%%%%%%%%%%%%%%%%%%%%%%%%%%%%%%%%%%%%%%%%%%%%%%%%%%%%
%%%%%%%%%%%%%%%%%%%%%%%%%%%%%%%%%%%%%%%%%%%%%%%%%%%%%%%
% %
% % The next command allows your in import encapsulated
% % postscript files, .epsf or .eps files, which
% % contain vector graphic image data.
% %
%%%%%%%%%%%%%%%%%%%%%%%%%%%%%%%%%%%%%%%%%%%%%%%%%%%%%%%
\usepackage{amssymb}
\usepackage{amsmath}
\usepackage{amsthm}
\usepackage{mathrsfs}
\usepackage{graphicx}
\usepackage[utf8]{inputenc}
\usepackage{multirow}
\usepackage{color}
%%%%%%%%%%%%%%%%%%%%%%%%%%%%%%%%%%%%%%%%%%%%%%%%%%%%%%%
% % We use newtheorem to define theorem-like structures
% %
% % Here are some common ones. . .
%%%%%%%%%%%%%%%%%%%%%%%%%%%%%%%%%%%%%%%%%%%%%%%%%%%%%%%
\newtheorem{theorem}{Theorem}
\newtheorem{lemma}{Lemma}
\newtheorem{proposition}{Proposition}
\newtheorem{scolium}{Scolium}   %% And a not so common one.
\newtheorem{definition}{Definition}
%\newenvironment{proof}{{\sc Proof:}}{~\hfill QED}
\newenvironment{AMS}{}{}
\newenvironment{keywords}{}{}
%%%%%%%%%%%%%%%%%%%%%%%%%%%%%%%%%%%%%%%%%%%%%%%%%%%%%%%
% %   The first thanks indicates your affiliation
% %
% %  Just the name here.
% %
% % Your mailing address goes at the end.
% %
% % \thanks is also how you indicate grant support
% %
%%%%%%%%%%%%%%%%%%%%%%%%%%%%%%%%%%%%%%%%%%%%%%%%%%%%%%%


\begin{document}
\newpage
\maketitle
%%%%%%%%%%%%%%%%%%%%%%%%%%%
% abstract, keywords and Subject classification are optional.
%%%%%%%%%%%%%%%%%%%%%%%%%%%
\begin{abstract}
    Let $D$ be a set of positive integers. A $D$-diffsequence of length $k$ is a sequence of positive integers $a_1 <  \cdots < a_k$ such that $a_{i+1}-a_i\in D$ for $i=1,\ldots,k-1$. For $D=\{2^i\mid i\in \mathbb{Z}_{\ge 0}\}$, it is known that there exists a minimum integer $n$, denoted by $\Delta(D,k)$, such that every $2$-coloring of $\{1,\ldots n \}$ admits a monochromatic $D$-diffsequence of length $k$. In this work, we prove a new lower bound for $\Delta(D,k)$ to $\Delta(D,k)\ge \left(\sqrt{\frac{8k-5}{12}}-\frac12\right)2^{\left(\sqrt{\frac{8k-5}{3}}-3\right)}$, asymptotically improving the exponential constant in the bound proved by Clifton.
\end{abstract}

% Most people don't use these, so they are "commented out"
% by starting the lines with a "%"
%\begin{keywords}
%   \LaTeX, typesetting
%\end{keywords}

%\begin{AMS}
%   50C60, 18C25
%\end{AMS}

%%%%%%%%%%%%%%%%%%%%%%
% % Here is the start of the Text
%%%%%%%%%%%%%%%%%%%%%%
\section{Introduction}
An $r$-coloring of a set $A$ is a function $\chi : A \to \{0, \ldots, r-1\}$, where each element of $A$ is assigned one of $r$ distinct colors. We say a subset $A' \subseteq A$ is monochromatic under $\chi$ if $\chi$ is constant over $A'$. In 1927, Van der Waerden~\cite{Waerdan} showed that for any $r$-coloring of the positive integers, at least one color is such that there are monochromatic arithmetic progressions of arbitrary lengths. Equivalently, for any positive integer $k$, there exists a minimum integer $n$, such that for any $r$-coloring of $\{1,\ldots, n \}$, there must be a monochromatic arithmetic progression of length $k$. Motivated by generalizing this property to other sequences, Landman and Robertson~\cite{LR} considered the following question: Does an $r$-coloring of positive integers produce a monochromatic sequence of arbitrary length, whose gaps lie in some fixed set $D$? In particular, they introduced the notion of a $D$-diffsequence.
\begin{definition}
    Given a set $D$ of positive integers, a $D$-\textit{diffsequence} of length $k$ is a sequence of positive integers $a_1 < \dots < a_k$ such that $a_{i+1}-a_i\in D$ for $i \in \{1,\dots,k-1\}$.
\end{definition}
\begin{definition}
    A set $D$ is called $r$-accessible if every $r$-coloring of the positive integers contains arbitrarily long monochromatic $D$-diffsequences.
\end{definition}
Several works have sought to characterize the accessibility of canonical integer sets $D$ (see \cite{book}), such as the set of Fibonacci numbers \cite{LR, Wesley}, primes \cite{LR}, and fixed translates of primes \cite{LV10}, and powers of integers. Consider the set of powers of some integer $t \in \mathbb{Z}_{> 0}$, i.e. $D = \{ t^{i}   \mid i \in \mathbb{Z}_{\ge 0}\}$. For $t\ge 3$, there is a periodic coloring modulo $t-1$ that avoids arbitrarily long $D$-diffsequences. It is known that $D$ is only $2$-accessible when $t = 2$ \cite{book}. Analogous to Van der Waerden's inquiry, the problem of $r$-accessibility of a set $D$ can be rephrased to finding the minimum integer $n$, denoted by $\Delta(D,k;r)$, such that every $r$-coloring of $\{1, \ldots, n\}$ admits a monochromatic $D$-diffsequence of length $k$. 

\begin{definition}
    Let $\Delta(D,k;r)$ be the smallest positive integer $n$, such that for any $r$-coloring of $\{1, \ldots, n \}$, we can find a $k$ length monochromatic $D$-diffsequence. Let $\Delta(D,k) = \Delta(D,k;2)$.
\end{definition}

For the set $D$ of powers of $2$, i.e., $D = \{ 2^i \mid i \in \mathbb{Z}_{\ge 0}\}$, Landman and Robertson \cite{LR} showed that $D$ is $2$-accessible, with $\Delta(D,k) \le 2^k - 1$. Chokshi et. al. \cite{CCLS18} conjectured that $\Delta(D,k)$ can be lower bounded exponentially. Clifton \cite{Clifton} established a lower bound which is asymptotically $\mathcal{O}(2^{\sqrt{2k}})$. In particular,
\begin{equation}
    \Delta(D, k) \geq 2^{\sqrt{2k}} \left( \frac{(\sqrt{2} - 1)k}{8} - \frac{\sqrt{k}}{8} \right) + \frac{\sqrt{k}}{2}.
\end{equation}
In this work, we establish a new lower bound for $\Delta(D,k)$, asymptotically improving the exponential constant in Clifton's bound from $\mathcal{O}(2^{\sqrt{2k}})$ to $\mathcal{O}(2^{\sqrt{8k/3}})$. Our main result is stated as follows.  
\begin{theorem}\label{thm: MainThm}
    For $D=\{2^i\mid i\in\mathbb{Z}_{\ge0}\}$ we have
    \begin{equation}\label{Theone}
        \Delta(D,k)\ge \left(\sqrt{\frac{8k-5}{12}}-\frac12\right)2^{\left(\sqrt{\frac{8k-5}{3}}-3\right)}.
    \end{equation}
\end{theorem}



% Clifton considers the coloring $P_t$ of the first $2^t$ positive integers represented by the first $2^t$ bits of the Thue-Morse sequence. They show that there are at most $m+1$ gaps of size $2^m$ for $m<t-1$ and there is no gap of size $2^{t-1}$. Summing this up gives a bound of $2^{\lceil\sqrt{2k}\rceil}\left(\frac{\sqrt{2k}}{2}-1\right)$. Clifton further improves this bound by a factor of $\Theta(\sqrt{k})$ though still keeping the exponent the same. \newline
% To prove this Clifton considers the coloring $P_{t,u}$ which is obtained by replacing each bit in $P_t$ by $2^u$ copies of itself, which they call a sub-block. They now bound the number of gaps of length $2^m$. For this, they consider a difsequence in $P_{t,u}$ and look at the different sub blocks the diffsequence contains. This also turns out to be a diffsequence in $P_t$. Clifton then uses the bounds shown previously, combining with how they affect the new structure, gives the lower bound on $\Delta(D,k)$. 

% We use a similar strategy of shrinking our string and looking at diffsequences in the smaller string and trying to bound their difference. Clifton's blocks contain the same color repeated $2^u$ times for some positive integer $u$. The blocks we use ($\mathbf{a}$ and $\mathbf{b}$, Definition \ref{find pos}) are quite different. We also do not directly bound the length of the monochromatic diffsequence of both strings; instead, we bound the length of the alternating diffsequence (Definition \ref{def: hops}). Further, we will repeat our contraction process multiple times to bound the diffsequence of our original string.



% =======================
% === Preliminaries ==== 
% =======================

\section{Construction}\label{sec: construction}
We will begin by formally constructing the coloring function $\kappa$ used to show the bound in Theorem~\ref{thm: MainThm}. We define a mapping from some coloring of $l$ integers to a coloring of $4l$ integers, which we call an \textit{expansion}. The process of \textit{expansion} is then recursively used to construct the required coloring $\kappa$. To prove the lower bound in Theorem~\ref{thm: MainThm}, we explicitly define a $2$-coloring scheme of $\{1,\ldots, l \cdot  4^{l-1}\}$ which avoids any monochromatic diffsequences of length $k = (3l^2-3l+2)/2$.  For a coloring $\chi \colon\{1,\ldots, l\} \to \{0,1\}$, define its \textit{expansion} $\chi^{(1)}\colon \{1,\ldots, 4l\} \to \{0,1\}$ by setting, for $i \in \{1, \ldots,  l\}$ and $j \in [0,3]$, 
\begin{equation}
\chi^{(1)}(4i - j) =
\begin{cases}
\chi(i), & \text{if } j \in \{2,3\}, \\
1 - \chi(i), & \text{if } j \in \{0,1\}.
\end{cases}
\end{equation}
The coloring scheme is constructed recursively using the process of expansion. Repeating this procedure $r$ times, we eventually obtain a coloring of the first $l \cdot 2^{2r}$ positive integers, which we call $\chi^{(r)}$. To show the lower bound in Theorem~\ref{thm: MainThm}, begin with the trivial coloring: $\chi_{\mathrm{alt}} \colon \{1,\ldots, l\} \to \{0,1\}$ such that $\chi_{\mathrm{alt}}(i) = i \mod{2}$. Then, we will show that the coloring $\chi_{\mathrm{alt}}^{(l-1)} \colon [ l \, \cdot \, 4^{l-1}] \to \{0,1\}$ avoids monochromatic diffsequences of length $k$, i.e. $\kappa \equiv  \chi_{\mathrm{alt}}^{(l-1)}$. 


\section{Preliminaries}\label{prev}
In this section, we will rephrase $2$-coloring functions, and the process of expansion to their more intuitive analogue, binary strings. For positive integers \( N \), by \(  [N] \) we mean the set \( \{1, 2, 3, \ldots, N\} \). By a 2-coloring of \([N] \) we mean a mapping \( \chi :[N] \to \{0, 1\} \). For the rest of this work, we will use the notation $D$ to denote the set of all non-negative integer powers of two, i.e., $D=\{2^i\mid i\in\mathbb{Z}_{\ge0}\}$.
Note that any $2$-coloring of $[n]$ can be represented as a binary string $\mathbf{s} = s_1 \ldots s_n$, where $s_i = \chi(i)$. For a string $\mathbf{s}$ corresponding to a given $2$-coloring of $[n]$, a monochromatic diffsequence $\mathcal{A} = \{a_i \}_{i=1}^{k}$ is such that $a_{i+1}-a_i \in D$ and $s_{a_{i+1}} = s_{a_i}$, i.e., $\chi(a_{i+1}) = \chi(a_i)$, for $i \in [k-1]$.
\begin{definition}\label{expand}
    Let $\mathbf{s}$ be a binary string of length $l$. The expansion of $\mathbf{s}$ is defined as a binary string $\mathbf{s}^{(1)}$ of length $4l$, obtained by replacing each $0$ in $\mathbf{s}$ by $0011$ and each $1$ in $\mathbf{s}$ by $1100$. The string $\mathbf{s}^{(r)}$ is recursively defined as the expansion of $\mathbf{s}^{(r-1)}$, where $r \ge 1$ and $\mathbf{s}^{(0)} = \mathbf{s}$.
\end{definition}

% \begin{definition}\label{cont def}
%     A binary string $\mathbf{s}'$ is contractable if there exists a binary string $\mathbf{s}$ such that $\mathbf{s}' = \mathbf{s}^{(1)}$.
% \end{definition}

\begin{ex}\label{bigex}
    Consider the string $\mathbf{t}=1011$, its expansion is $\mathbf{t}^{(1)} = \overbrace{1100}_{\phantom{.}}\,\overbrace{0011}_{\phantom{.}}\,\overbrace{1100}_{\phantom{.}}\,\overbrace{1100}_{\phantom{.}}$. 
\end{ex}
In Lemma~\ref{MainLemma}, we will bound the length of certain diffsequences of an expanded string. Therefore, we wish to understand how diffsequences behave under expansion. Let $\mathbf{s}^{(1)}$ be the expansion of some string $\mathbf{s}$. The process of expansion is such that the $i^{th}$ bit of $\mathbf{s}^{(1)}$ is determined uniquely by the $\lceil i/4 \rceil^{th}$ bit of the string $\mathbf{s}$. This lets us determine the \textit{block}, either $0011$ or $1100$, to which the $i^{th}$ bit of $\mathbf{s}^{(1)}$ belongs. For convenience, we call this mapping the position of $i$.
\begin{definition}\label{pos}
    Let the position of $i \in \mathbb{Z}_{+}$ be defined as the function $\text{pos}(i)=\left\lceil i/4 \right\rceil$.
\end{definition}

% \begin{ex}\label{findex}
%     Consider the $16$ length string $\mathbf{t}^{(1)}=1100001100111100$ from Example~\ref{hop example}. The contraction of $\mathbf{t}^{(1)}$ is the $4$ length binary string $t:=1001$. Therefore we have $\text{find}(5)=A$ as $t_{\text{pos(5)}}=t_2=0$.
% \end{ex}

The choice of $0011$ and $1100$ as the types of blocks used in expansion is motivated by the following two reasons: a) The length of blocks being $4$ implies that for any two bit indices $i$ and $i'$ in the expanded string $\mathbf{s}^{(1)}$ that differ by a power of $2$, i.e. $i'-i \in D$, their corresponding positions in the string $\mathbf{s}$ are such that $\mathrm{pos}(i') - \mathrm{pos}(i) \in D \, \cup \, \{0\}$; and b) For a monochromatic diffsequence $a_1< \ldots < a_k$ in the expanded string which is such that some of its consecutive elements lie in different blocks i.e. $s_{pos(a_{j+1})} \neq s_{pos(a_{j})}$, for some $j \in [k-1]$, we must have that $a_{j+1} - a_{j} \le 2$, i.e. they must lie in consecutive blocks. This is enabled by the fact that the blocks $0011$ and $1100$ are conjugates of each other. This will let us reduce monochromatic diffsequences in the expanded string $\mathbf{s}^{(1)}$ to two contiguous diffsequences of opposite colors in $\mathbf{s}$, partitioned by a $\textit{hop}$. A hop is defined to be an index in a diffsequence, at which the assigned color of the diffsequence changes.

\begin{definition}\label{def: hops}
Let $\mathbf{s} \in \{ 0,1 \}^n$ and $\mathcal{A}=\{a_i\}_{i=1}^k$ a $D$-diffsequence. A $\textit{hop}$ is defined as an index $j \in [k-1]$ such that $a_{j+1}-a_j=1$ and $s_{a_{j+1}}\neq s_{a_{j}}$. Define $H_{\mathcal{A}}$ to be the set of all hops in $\mathcal{A}$.
\end{definition}

\begin{definition}\label{def: Psi_psi}
For a string $\mathbf{s} \in \{ 0,1 \}^n$, define $\Psi_{\mathbf{s}}(h)$ to be the set of all diffsequences with at most $h$ hops. Formally, $ \Psi_{\mathbf{s}}(h): =\{\mathcal{A}=\{a_i\}_{i=1}^k 
    :\mathcal{A} \text{ is a } D\text{-diffsequence with } |H_{\mathcal{A}}|\le h\}$. Let $\psi_\mathbf{s}(h)$ denote the maximum length of any diffsequence in $\Psi_{\mathbf{s}}(h)$, and $\psi_\mathbf{s}(h) = 0$ if $\Psi_{\mathbf{s}}(h) = \emptyset$.
\end{definition}

\begin{rem}
    $\Psi_{\mathbf{s}}(0)$ is the set of all monochromatic diffsequences of $\mathbf{s}$. Therefore, if the length of $\mathbf{s}$ is greater than $\Delta(D,k)$, the maximum length of elements of $\Psi_{\mathbf{s}}(0)$ is at least $k$.
\end{rem}

\begin{ex}\label{hop example}
Consider the string $\mathbf{s}=1100001100111100$ of length $16$. The sequence of integers $\mathcal{A} = \{1,2,3,5,9\}$ is an example of a $D$-diffsequence. The corresponding subsequence of $\mathbf{s}$ is $s_1s_2s_3s_5s_9 = 11000$, and $H_{\mathcal{A}} = \{2\}$. Since this is the only hop in the diffsequence, $\mathcal{A} \in \Psi_{\mathbf{s}}(1)$.   
\end{ex}

%        Then the corresponding $2$ coloring of this diffsequence is represented by arrows on the positions in 
%        $\mathbf{s}$ \[\overset{\downarrow}{1}\overset{\downarrow}{1}\overset{\downarrow}{0}0\overset{\downarrow}{0}011\overset{\downarrow}{0}0111100 = 11000.\]

% \begin{definition}\label{find pos}
%     Consider the binary string $\mathbf{s}^{(1)}$ of length $4l$, the expansion of some string $\mathbf{s}$. Let $\mathbf{a}=0011$ and $\mathbf{b}=1100$.  Define a function $\text{find}:\{0,\dots,4n-1\}\to\{A,B\}$ such that 
%     \begin{equation}
%         \text{find}(i)=
%         \begin{cases}
%         A\quad \text{if $\mathbf{t}_{\text{pos}(i)}=0$},\\
%         B\quad \text{if $\mathbf{t}_{\text{pos}(i)}=1$}.
%         \end{cases}
%     \end{equation}
% \end{definition}

% Thus, $\textit{find}$ tells us in which block ($0011$ or $1100$) a position in $\mathbf{s}$ lies. We restrict the output of $\textit{find}$ to two symbols $A$ and $B$ which correspond to the strings $0011$ and $1100$. Notice that $\mathbf{a}=0011$ and $\mathbf{b}=1100$ are the same strings that were used in expansion and contraction (Definition~\ref{expand} and~\ref{cont def})


 
% Our aim is going to be to convert a diffsequence $\{a_i\}_{i=1}^k$ in $\mathbf{s}$ into a diffsequence in $\mathbf{t}$, by contracting each $a_i$ into the block it correspond in $\mathbf{t}$, which formally is $\lceil\frac{a_i}{4}\rceil=\text{pos}(a_i)$.

% \begin{rem}
%     The idea of looking at a diffsequence and converting it to another diffsequence in a smaller string is similar to the proof in \cite{Clifton}.
% \end{rem}

% which states that if the color of our diffsequence is $1$ then we can never go from a $B$ to an $A$ but, we can go from $A$ to $B$, though only if the $A$ and $B$ are consecutive and we cannot do this hop anymore

% In Proposition~\ref{x-y prop}, 
% we show how monochromatic diffsequences behave under contraction. In particular, given a binary string $\mathbf{s}$, we examine the values of $\textit{find}$ and $\textit{pos}$ for two consecutive elements of the same color in a diffsequence, i.e., the positions of these two elements differ by a power of $2$. 

%We will also show how diffsequences are almost preserved under contraction in the following proposition.

%The following proposition tells some properties about the diffsequences in contracted coloring. This will be the key to proving the base case of Lemma~\ref{MainLemma}. This proposition tells properties about $\textit{find}$ and $\textit{pos}$ for two numbers differing by a power of two.
%\[\tikzmarknode{aloo}{1}0010\tikzmarknode{baingan}{1}101\tikzmarknode{gobi}{1}003\]

%\begin{tikzpicture}[remember picture,overlay]
%\draw[-latex]([yshift=.5ex]aloo.north) to[bend left]node[above]{} %([yshift=.5ex]baingan.north);
%\draw[-latex]([yshift=.5ex]baingan.north) to[bend left]node[above]{} ([yshift=.5ex]gobi.north);
%\end{tikzpicture}
The following proposition and its corollaries give us a recipe to reduce diffsequences with $h$ hops in some expanded string $\mathbf{s}^{(1)}$, to diffsequences with $h+1$ hops in $\mathbf{s}$. Proposition~\ref{prop: reduce_diffseq} shows that in the expanded string $\mathbf{s}^{(1)}$, if two positions $\{x,y\}$ of some diffsequence have the same color $c \in \{0, 1 \}$, and their corresponding positions $\{\text{pos}(x), \text{pos}(y) \}$ in $\mathbf{s}$ have different colors, then $\text{pos}(y)=\text{pos}(x)+1$ i.e. $x$ and $y$ lie in blocks next to each other. Moreover, we must have $s_{\text{pos}(x)}=1-c$, and $ s_{\text{pos}(y)}=c$. 

\begin{proposition}\label{prop: reduce_diffseq}
    For a binary string $\mathbf{s}$,
    let $x,y$ be positive integers such that $y-x\in D$ and $s^{(1)}_x=s^{(1)}_y$. If $s_{\text{pos}(x)}\neq s_{\text{pos}(y)}$, we must have $s_{\text{pos}(x)}=1-s^{(1)}_x$, $ s_{\text{pos}(y)}=s^{(1)}_y$, and $\text{pos}(y)=\text{pos}(x)+1$.
\end{proposition}

\begin{proof}
The condition $s_{\text{pos}(x)}\neq s_{\text{pos}(y)}$ means that $s_x^{(1)}$ and $s_y^{(1)}$ belong to different types of blocks, $0011$ and $1100$. Without loss of generality, assume that $s_x^{(1)} = s_y^{(1)}$ = 1. The proof is divided into two cases according to the value of $y-x$. (a) For $y-x \in \{1,2 \}$, we have $\text{pos}(y)-\text{pos}(x)\in \{0, 1\}$ which means that $s_x^{(1)}$ and $s_y^{(1)}$ belong either to the same block or consecutive blocks. However, since $s_x^{(1)}$ and $s_y^{(1)}$ must belong to different types of blocks, these must be consecutive.
Therefore $\text{pos}(y)-\text{pos}(x)=1$. The condition $s_x^{(1)} = s_y^{(1)} = 1$ is only possible when 
$s_x^{(1)} \in \{0011\}$, and $s_y^{(1)} \in \{1100\}$, which implies $s_{\text{pos}(x)}=0$ and $s_{\text{pos}(y)}=1$. (b) For $y-x=2^m$, $x \equiv y \mod{4}$ and therefore, $s_x^{(1)}$ and $s_y^{(1)}$ occupy the same position within their respective blocks. Since the two types of blocks, $0011$ and $1100$ are conjugates of each other, $s^{(1)}_{x} = s^{(1)}_{y}$ is only possible when the blocks are identical. This would mean that $s_{\text{pos}(x)}$ and $s_{\text{pos}(y)}$ are identical, which is a contradiction.
\end{proof}

% The following corollary is obtained by taking conjugates in Proposition~\ref{x-y prop}, i.e., replacing all $0$'s with $1$'s and all $1$'s with $0$'s.
 
% \begin{cor}
%     Let $\mathbf{s}$ be any contractable binary string. Let $x,y$ be positive integers such that $x-y \in D$, and $s_x=s_y=0$, then
%     \begin{itemize}
%         \item $\text{pos}(x)-\text{pos}(y) \in D \cup \{ 0 \}$.
%         \item If $\text{find}(y)=A$, then $\text{find}(x)=A$.
%         \item If $\text{find}(y)=B$ and $\text{find}(x)=A$, then $\text{pos}(x)=\text{pos}(y)+1$.
%     \end{itemize}
% \end{cor}

% We now prove some facts about how $\textit{find}$ and $\textit{pos}$ act on a particular diffsequence with some constraint.  In corollaries~\ref{A B split cor} and~\ref{pos bound cor} we only consider the case in which the color of the diffsequence is always $1$. The other case, when the color is $0$, can be obtained by taking conjugates. 

Let $x_1<\dots<x_k$ be a monochromatic diffsequence of color $c \in \{0,1\}$ in the string $\mathbf{s}^{(1)}$, which is the expansion of some string $\mathbf{s}$. As noted earlier, $y-x \in D$, implies $\text{pos}(y)-\text{pos}(x) \in D \ \cup \{0\}$; the sequence characterized by the set $\{\text{pos}(x_i)\}_{i=1}^{k}$ is a $D$-diffsquence in $\mathbf{s}$. A consequence of Proposition~\ref{prop: reduce_diffseq} is that the substring $s_{\text{pos}(x_i)}$ can be split into two contiguous monochromatic substrings, the first with the color $1-c$ and the other with $c$. This result is proved formally in Corollary~\ref{A B split cor}. Therefore, a monochromatic diffsequence in $\mathbf{s}^{(1)}$ can be mapped to a diffsequence with one switch in color in $\mathbf{s}$, i.e., with one hop (Definition~\ref{def: hops}). This is a mapping between elements of $\Psi_{\mathbf{s^{(1)}}}(0)$ to elements of $\Psi_{\mathbf{s}}(1)$. In Corollary~\ref{pos bound cor}, we will establish that the set $\{\text{pos}(x_i)\}_{i=1}^{k}$ has at least $k-1$ distinct elements under certain conditions. These two corollaries together help in bounding $\psi_{\mathbf{s^{(1)}}}(0)$ in terms of $\psi_{\mathbf{s}}(1)$ (Definition~\ref{def: Psi_psi}) as $\psi_{\mathbf{s^{(1)}}}(0) \le \psi_{\mathbf{s}}(1) +2$.

\begin{cor}\label{A B split cor}
    Let $\mathbf{s}^{(1)}$ be an expansion of a binary string $\mathbf{s}$ and $x_1<\dots<x_k$ be a monochromatic $D$-diffsequence in $\mathbf{s}^{(1)}$ with color $c \in \{0,1\}$. Then, there exists a non negative integer $m\le k$ such that $s_{\text{pos}(x_j)}=1-c$ if and only if $j\le m$.
\end{cor}
\begin{proof}
    From Proposition~\ref{prop: reduce_diffseq} we have the following possibilities for each $i\le k-1$: $s_{\text{pos}(x_i)}=s_{\text{pos}(x_{i+1})}$ or $s_{\text{pos}(x_i)}=1-c$ and $s_{\text{pos}(x_{i+1})}=c$. Now, if there exists an index $j$ for which $s_{\text{pos}(x_j)}=1-c$ and $s_{\text{pos}(x_{j+1})}=c$ (if not, then we are already done), then pick the smallest such $j$. Thus, $s_{\text{pos}(x_i)}=1-c$ for all $i\le j$. We must have $s_{\text{pos}(x_i)}=c$ for all $i>j$, as otherwise there would exist an index $j^{\prime}$ such that $s_{\text{pos}(x_{j^{\prime}})}=c$ and $s_{\text{pos}(x_{j^{\prime}+1})}=1-c$, which is not possible. Hence, proving our claim.
\end{proof}

\begin{cor}\label{pos bound cor}
   Let $\mathbf{s}^{(1)}$ be an expansion of a binary string $\mathbf{s}$ and $x_1<\dots<x_k$ be a monochromatic $D$-diffsequence in $\mathbf{s}^{(1)}$ with color $c\in\{0,1\}$. Suppose that there exists $d\in\{0,1\}$ such that $s_{\text{pos}(x_i)}=d$ for all $1\le i\le k$. Then the set $\{\text{pos}(x_i)\}^k_{i=1}$ has at least $k-1$ distinct elements.
\end{cor}
\begin{proof}
Without loss of generality, assume $c=1$. The proofs for $d=0$ and $d=1$ are analogous; therefore, we will further assume $ d=0$. Since $\{\text{pos}(x_i)\}^k_{i=1}$ is a non-decreasing sequence, to prove it has at least $k-1$ distinct elements, it suffices to show that $\text{pos}(x_j)=\text{pos}(x_{j+1})$ is true for at most one index. If it exists, let $j$ be such an index. For each $1\le i\le k$, the conditions $s^{(1)}_{x_i}=1$ and $s_{\text{pos}(x_i)}=0$ imply that $s^{(1)}_{x_i}$ corresponds to the third or the fourth position in the block $0011$. Furthermore, if $s^{(1)}_{x_i}$ corresponds to the fourth position in the block $0011$, then $s^{(1)}_{x_{i+1}}$ must also correspond to the fourth position. This is because if it corresponds to the third position, then the difference between $x_{i+1}$ and $x_i$ would be odd, implying it to be $1$, but this is clearly not possible. For an index $j$ to satisfy $\text{pos}(x_j)=\text{pos}(x_{j+1})$, we must have $s^{(1)}_{x_j}$ and $s^{(1)}_{x_{j+1}}$ correspond to the third and fourth position in the block $0011$ respectively. This means that all indices at least $j+1$ must correspond to the fourth position in the block, thus enforcing the impossibility of another index $i$ satisfying $\text{pos}(x_i)=\text{pos}(x_{i+1})$.
\end{proof}

%In SMTH it is proved $\Delta(D,k)\leq 2^k-1$ and SMTH proved \begin{equation}\Delta(D,k)\ge2^{\sqrt{2k}}\left(\frac{(\sqrt{2}-1)k}{8}-\frac{\sqrt{k}}{8}\right)+\frac{\sqrt{k}}{2}\end{equation} we improve this bound.

Corollaries \ref{A B split cor} and \ref{pos bound cor} show that the process of expansion naturally induces a mapping from monochromatic diffsequences in an expanded string and diffsequences with one hop in the original string. Motivated by this approach, in Lemma~\ref{MainLemma}, we inductively extend this reasoning to obtain a mapping from elements in $\Psi_{\mathbf{s}^{(1)}}(h)$ to $\Psi_{\mathbf{s}}(h+1)$ for some string $\mathbf{s}$. This will help us obtain a bound on the maximum length of a monochromatic diffsequence in our explicit construction.

%LEQ OR = !!!!!!!!!!!!!!!!!!!!!!!!!!!!!!!!!!




\section{Main result}\label{MainResult}
To prove the main result of this work, Theorem~\ref{thm: MainThm}, we explicitly need to show that the construction in Section~\ref{sec: construction} does not admit any monochromatic $D$-diffsequences of length $k$. For any string $\mathbf{s}$ and its expansion $\mathbf{s}^{(1)}$, Lemma~\ref{MainLemma} gives a bound on the difference between the maximum length of a diffsequence with at most $h+1$ hops in $\mathbf{s}$, $\psi_{\mathbf{s}}(h+1)$, and the
maximum length of a diffsequence with $h$ hops in the expanded string $\mathbf{s}^{(1)}$, $\psi_{\mathbf{s}^{(1)}}(h)$ (Definition~\ref{def: Psi_psi}). Formally, for any binary string $\mathbf{s}$, Lemma~\ref{MainLemma} shows that $\psi_{\mathbf{s}^{(1)}}(h) \le \psi_{\mathbf{s}}(h+1) + 3h + 2$. Recall that in our explicit construction, we begin with the initial coloring $\chi_{\text{alt}} \colon [l] \to \{0,1\}$ such that $\chi_\text{alt}(i) = i \mod{2}$. This corresponds to the binary string $\mathbf{s} = 1010 \ldots$ of length $l$. This string is such that the diffsequence $\{i\}_{i=1}^{l}$ has $l-1$ hops, i.e. $\{i\}_{i=1}^{l} \in \Psi_{\mathbf{s}}(l-1)$ and $\psi_{\mathbf{s}}(l-1) = l$. Therefore, using Lemma~\ref{MainLemma} recursively, we can show that $\psi_{\mathbf{s}^{(l-1)}}(0) \le (3l^2-3l+2)/2$. That is, we show that the maximum possible length of a monochromatic diffsequence in $\mathbf{s}^{(l-1)}$ is at most $(3l^2-3l+2)/2$. The proof of Lemma~\ref{MainLemma} uses induction on the number of hops, exploiting the fact that diffsequences with $h$ hops, i.e., diffsequences in $\Psi_{\mathbf{s}^{(1)}}(h)$, are a concatenation of diffsequences with $h-1$ hops and monochromatic diffsequences. For an expanded string, we can therefore inductively apply Corollaries \ref{A B split cor} and \ref{pos bound cor}.

% we contract our given string, and then we look at how the diffsequence in the original string corresponds to the contracted string. Thus, we are looking at the diffsequence as hops from ones and zeros, whereas in the contracted string, we look at it as hops from $A$'s and $B$'s. The first part of the proof presents the base case of $h=0$. The main idea for this case is the second property of Proposition~\ref{x-y prop} which states that if the color of our diffsequence is $1$ then we can never go from a $B$ to an $A$ but, we can go from $A$ to $B$, though only if the $A$ and $B$ are consecutive and we cannot do this hop anymore, this explains the increase of hops when we are contracting. The second part tries to find the first time we hop, then splits the string into two parts, the first having one or zero hops and the second having the remaining $h-1$. We then apply the induction hypothesis to the remaining string to bound the length of the contracted string. The following proves the lemma in a much more formal way.

% \begin{ex}
%     For example, we look at the diffsequence with one hop from Example~\ref{hop example}. The $16$ length string $\mathbf{s}=s_1s_2\dots s_{16}=1100001100111100$. We examine the diffsequence $1,2,3,5,9$.  We represent this by arrows over the positions in $\mathbf{s}$ \[\overset{\downarrow}{1}\overset{\downarrow}{1}\overset{\downarrow}{0}00\overset{\downarrow}{0}110\overset{\downarrow}{0}111100.\]

%     Now the contraction of $\mathbf{s}$ is the string $1001$. The above diffsequence corresponds to a new diffsequence in $1001$ which is $1,2,3$ as shown below.
%     \[\overset{\downarrow}{1}\overset{\downarrow}{0}\overset{\downarrow}{0}1.\]
% \end{ex}

\begin{lemma}\label{MainLemma}
Let $\mathbf{s}$ be a binary string and let $\mathbf{s}^{(1)}$ denote the expansion of $\mathbf{s}$. Then 
\begin{equation}
    \psi_{\mathbf{s}^{(1)}}(h) \le \psi_{\mathbf{s}}(h+1) + 3h + 2. 
\end{equation}

\end{lemma}


\begin{proof}
Using induction on $h$, we will show that for any diffsequence in $\Psi_{\mathbf{s^{(1)}}}(h)$ of length $\ell$, there is a corresponding diffsequence in $\Psi_{\mathbf{s}}(h+1)$ with length at least $\ell-3h-2$. To prove the base case, $h=0$, let $ \mathcal{P}_0 =\{ p_i\}_{i=1}^{\ell} \in \Psi_{\mathbf{s}^{(1)}}(0)$ be a monochromatic diffsequence in $\mathbf{s}^{(1)}$ of length $\ell$. Without loss of generality, assume that $ \mathcal{P}_0$ has color $c=1$ in $\mathbf{s}^{(1)}$, i.e., $s^{(1)}_{p_i} = 1$, for $1 \le i \le \ell$. Now, consider the diffsequence $\text{pos}(\mathcal{P}_0) = \{\text{pos}(p_i)\}_{i=1}^{\ell}$ in $\mathbf{s}$. By Corollary~\ref{A B split cor}, we know that the substring $\{s_{\text{pos}(x_i)}\}_{i=1}^{\ell}$ can be split into at most two contiguous monochromatic substrings, the first with the color $1-c=0$ and the other with $c=1$. Let $\mathcal{P}_{\text{left}}$ and $\mathcal{P}_{\text{right}}$ be (possibly empty) subsequences of $ \mathcal{P}_0$ such that $\text{pos}(\mathcal{P}_{\text{left}})$ has color $0$ in $\mathbf{s}$, and $\text{pos}(\mathcal{P}_{\text{right}})$ has color $1$ in $\mathbf{s}$. Using Corollary~\ref{pos bound cor}, note that $|\text{pos}(\mathcal{P}_{\text{left}})|\ge|\mathcal{P}_{\text{left}}|-1$ and $|\text{pos}(\mathcal{P}_{\text{right}})|\ge|\mathcal{P}_{\text{right}}|-1$. Because $\text{pos}(\mathcal{P}_0)$ is the union of $\text{pos}(\mathcal{P}_{\text{left}})$ and $\text{pos}(\mathcal{P}_{\text{right}})$, and both $\text{pos}(\mathcal{P}_{\text{left}})$ and $\text{pos}(\mathcal{P}_{\text{right}})$ are disjoint, we get that $|\text{pos}(\mathcal{P}_0)|=|\text{pos}(\mathcal{P}_{\text{left}})|+|\text{pos}(\mathcal{P}_{\text{right}})|\ge |\mathcal{P}_{\text{left}}|-1+|\mathcal{P}_{\text{left}}|-1=|\mathcal{P}|-2$.

To prove the inductive hypothesis, assume that the result is true for all $k\leq h-1$. Let $\mathcal{P}_h=\{p_i\}^\ell_{i=1}$ be a diffsequence in $\mathbf{s}^{(1)}$ with at most $h$ hops. If $\mathcal{P}_h$ has at most $h-1$ hops, i.e. $\mathcal{P}_h \in \Psi_{\mathbf{s}^{(1)}}(h-1)$, then using the inductive hypothesis for $k = h-1$, we are done. Now, let us consider the case where $\mathcal{P}_h$ has exactly $h \ge 1$ hops. Let $p_i\to p_{i+1}$ be the first time there is a hop in $\mathbf{s}^{(1)}$. Let $\mathcal{P}_{\text{left}}$ be the first $i$ elements of $\mathcal{P}_h$ and $\mathcal{P}_{\text{right}}$ be the remaining part. Analogous to the proof of the base case, we split $\mathcal{P}_h$ into two contiguous subsequences $\mathcal{P}_{\text{left}}$ and $\mathcal{P}_{\text{right}}$ such that the hop between $\mathcal{P}_{\text{left}}$ and $\mathcal{P}_{\text{right}}$ is the first hop in $\mathcal{P}_h$. This means that $\mathcal{P}_{\text{left}}$ is monochromatic in $\mathbf{s}^{(1)}$ and $\mathcal{P}_{\text{right}}$ has at most $h-1$ hops in $\mathbf{s}^{(1)}$. Notice that $\text{pos}(\mathcal{P}_h)$ is the union of $\text{pos}(\mathcal{P}_{\text{left}})$ and $\text{pos}(\mathcal{P}_{\text{right}})$. Furthermore, $\text{pos}(\mathcal{P}_{\text{left}})$ and $\text{pos}(\mathcal{P}_{\text{right}})$ intersect in at most one element. Hence, we can lower bound $|\mathcal{P}_h|$ by $|\text{pos}(\mathcal{P}_{\text{left}})|+|\text{pos}(\mathcal{P}_{\text{right}})|-1$. As $\mathcal{P}_{\text{left}}\in \Psi_{\mathbf{s}^{(1)}}(0)$ and $\mathcal{P}_{\text{right}}\in \Psi_{\mathbf{s}^{(1)}}(h-1)$, we can apply the induction hypothesis to both $\mathcal{P}_{\text{left}}$ and $\mathcal{P}_{\text{right}}$ to obtain the following. $|\text{pos}(\mathcal{P}_h)| \ge |\text{pos}(\mathcal{P}_{\text{left}})|+|\text{pos}(\mathcal{P}_{\text{right}})|-1 \ge |(\mathcal{P}_{\text{left}})|-2+|(\mathcal{P}_{\text{right}})|-3(h-1)-2-1 \ge |\mathcal{P}_h|-3h-2.$
\end{proof}
\noindent
\textbf{Proof of Theorem 1.} 
We now explicitly construct a coloring of length $l \cdot4^{l-1}$ with the length of the maximum diffsequence at most $(3l^2-3l+2)/2$. We start with the binary string $\mathbf{s}$ of length $l$ consisting of alternating ones and zeros, i.e, the one corresponding to the coloring $\chi(i)\equiv i\pmod{2}$. This ensures $\psi_{\mathbf{s}}(l-1)=l$. By Lemma~\ref{MainLemma}, $\psi_{\mathbf{s}^{(i)}}(l-1-i)\le \psi_{\mathbf{s}^{(i-1)}}(l-1-i+1)+3(l-i-1)+2$ for all $1\le i\le l-1$.
So, summing over all $1\le i\le l-1$ we get
\begin{equation*}
    \psi_{\mathbf{s}^{(l-1)}}(0)\le l+\frac{3(l-1)(l-2)}{2}+2(l-1)=\frac{3l^2-3l+2}{2}.
\end{equation*}
For any positive integer $k>\frac{3l^2-3l+2}{2}$, we obtain a coloring of length $l\cdot 4^{l-1}$ that avoids a monochromatic diffsequence of length $k$, therefore, implying that $\Delta(D,k)\ge l\cdot4^{l-1}$. Hence, choosing the largest such $l=\left\lfloor\sqrt{\frac{2(k-1)}{3}+\frac{1}{4}}+\frac{1}{2}\right\rfloor$, we obtain $\Delta(D,k)\ge\left(\sqrt{\frac{8k-5}{12}}-\frac12\right)2^{\left(\sqrt{\frac{8k-5}{3}}-3\right)}$.

%Hence, for any $k$ choosing maximal $l=\left\lfloor\sqrt{\frac{2(k-1)}{3}+\frac{1}{4}}+\frac{1}{2}\right\rfloor$. We get a $l  \cdot  4^{l-1}=l \cdot  4^{l-1}$ length sequence with no $k$ length diffsequence. Giving a lower bound of\[\left(\sqrt{\frac{8k-5}{12}}-\frac12\right)2^{\left(\sqrt{\frac{8k-5}{3}}-3\right)}.\]

%We have the following possibilities for a hop in $\mathbf{s}^{(1)}$: 
%\begin{itemize}
%    \item A hop from the last element of an $0011$ to the first element of a $1100$, where both these blocks are consecutive (call this the first case). Similarly, one from the last element from $1100$ to the first element of an $0011$. These look like $001\overset{\downarrow}{1}\overset{\downarrow}{0}011$ (and its conjugate).
%    \item A hop from the middle two elements of an $1100$ (call this the second case) and a similar one in $0011$. These look like $1\overset{\downarrow}{1}\overset{\downarrow}{0}0$ (and its conjugate).
%\end{itemize}





%------------------------------------------------------------\\

%We claim that in each case the number of times we switch in $S$ which is not counted in $T$ is %atmost 3 $\\$ 
%Define the set $X$ and the function $f$ as before. $\\$ Let $s_i\to s_{i+1}$ be the first time we hop in $S\\$
%If this falls in the first case, then we have that $\lceil\frac{s_{i}}{4}\rceil=\lceil\frac{s_{i+1}}{4}\rceil$, and we may also have that $\\\lceil\frac{s_{i+1}}{4}\rceil=\lceil\frac{s_{i+2}}{4}\rceil$ if $s_{i+2}=s_{i+1}+1$. Now let $i_2\to i_2+1$ be the next time we hop in $S$. Note that for all $i\leq j\leq i_2-1$, $f(s_j)=B$ and $\lceil\frac{s_{j}}{4}\rceil\neq\lceil\frac{s_{j+1}}{4}\rceil$.Furthermore there is at most one $j\leq i-1$ with $\lceil\frac{s_{j}}{4}\rceil=\lceil\frac{s_{j+1}}{4}\rceil$ (due to a similar proof as in the base case). By induction hypothesis on the sub string $4\lceil\frac{s_{i_2}}{4}\rceil$ to $|S|$ we have that this maps to a diffsequence (in $T$) with $l-1+1=l$ hops with length at least $u-i_2+1-3(l-1)-2$ , combining this with the diffsequence with 1 consecutive hop (in $T$) mapped by the diffsequence in the sub string $0$ to $4\lceil\frac{s_{i_2}}{4}\rceil-1 (\text{in } S)$   which has is of length at least $i_2-1-3$  (the minus three comes from the fact that there were atmost 3 indices $1\leq j\leq i_2-1$ with $\lceil\frac{s_{j}}{4}\rceil=\lceil\frac{s_{j+1}}{4}\rceil$) totally giving a atleast $u-3l-2$ length diffsequence with $l+1$ hops in $T\\\\$
%\noindent
%If this falls in the second case, Note that there is index $j\leq i-1$ such that $f(s_j)\neq f(s_{j+1})$ as otherwise we must have $s_j\equiv 3\pmod{4}$ and $s_{j+1}\equiv 0\pmod{4}$ henceforth for all $j+1\leq k\leq i$ we have $s_k\equiv1\pmod 4$ or $s_k\equiv 0\pmod 4$ but $s_i\equiv 3\pmod{4}$ a contradiction.
%And there is at most 1 index $j$ with $\lceil\frac{s_{j}}{4}\rceil=\lceil\frac{s_{j+1}}{4}\rceil$.  Therefore this diffsequence maps to a diffsequence (with no consecutive hops) in $T$ of length at least $i-1-1\\$
%Now we have $\lceil\frac{s_{i}}{4}\rceil\neq\lceil\frac{s_{i+1}}{4}\rceil$ and we may have $\lceil\frac{s_{i+1}}{4}\rceil=\lceil\frac{s_{i+2}}{4}\rceil$ if $s_{i+2}=s_{i+1}+1$, again as before define $i_2$ to be the next time we hop. Note that there is no $j$ with $\lceil\frac{s_{j}}{4}\rceil\neq\lceil\frac{s_{j+1}}{4}\rceil.\\$
%By induction hypothesis on the sub string $4\lceil\frac{s_{i_2}}{4}\rceil$ to $|S|$ we have that this maps to a diffsequence (in $T$) with $l-1+1=l$ hops with length at least $u-i_2+1-3(l-1)-2$ , combining this with the diffsequence with 1 consecutive hop (in $T$) mapped by the diffsequence in the sub string $0$ to $4\lceil\frac{s_{i_2}}{4}\rceil-1 (\text{in } S)$   which has is of length at least $i_2-1-2$  (the minus three comes from the fact that there were atmost 2 indices $1\leq j\leq i_2-1$ with $\lceil\frac{s_{j}}{4}\rceil=\lceil\frac{s_{j+1}}{4}\rceil$) totally giving a atleast $u-3l-1$ length diffsequence with $l+1$ hops in $T\\$
%and we are done.\\
%----------------------------------------------------------------------\\




%\noindent
%Consider the first time there is a hop in $\mathbf{s}^{(1)}$. This hop must either be the ones described in the first case or the second case, as we started with the color $1$ and never changed it before the first hop.
%Let $p_i\to p_{i+1}$ be the first time we hop.
%We split the proof into two cases depending on whether the hop $p_i\to p_{i+1}$ lies in the first case or in the second case.\\
%If the hop is like the one described in the first case.

%We split $\mathbf{s}^{(1)}$ into two smaller stings $\textbf{s}^{(1)}_1,\textbf{s}^{(1)}_2$ which are the expansion of some strings $s_1$ and $s_2$ respectively, as follows: Let $\textbf{s}^{(1)}_1$ be the first $p_i$ elements of $\mathbf{s}^{(1)}$, and $\textbf{s}^{(1)}_2$ be the remaining part. Both $\textbf{s}^{(1)}_1$ and $\textbf{s}^{(1)}_2$ are expansions of some strings $s_1$ and $s_2$, as $p_i\equiv0\pmod{4}$.\\
%The sequence $p_1,\dots,p_i$ is a diffsequence with $0$ hops in $\textbf{s}^{(1)}_1$, and $\mathbf{s_1}^{(1)}_{\text{pos}(p_i)}=0$ hence the diffsequnce this corresponds in $s_1$ must have no hops. This is because a hop would mean a change in color in $s_1$ which only happens when there exists a $j$ such that $\mathbf{s_1}^{(1)}_{\text{pos}(p_j)}\neq\mathbf{s_1}^{(1)}_{\text{pos}(p_{j+1})}$ which in this case isn't true. Notice this diffsequence satisfies the conditions of Corollary~\ref{pos bound cor}. Therefore, the diffsequence corresponding in $s_1$ must have length at least $i-1$ by the corollary. The induction hypothesis implies the diffsequence corresponding to $\{p_j\}_{j>i}$ in $s_2$ has length at least $\ell-i-3(h-1)-2$\\
%, and of length at least $i-1$. The sequence $p_{i+1},\dots p_{\ell}$ is a diffsequence with $h-1$ hops in $s_2$, so it corresponds to a diffsequence with at most $h$ hops in $t_2$ and of length at least $\ell-i-3(h-1)-2$ by induction hypothesis\\
%Adding the above results and the hop from $p_i\to p_{i+1}$ (which also is a hop in $\mathbf{s}$), we get a diffsequence with at most $h+1$ hops of length at least $1+\ell-i-3(h-1)-2+i-1=\ell-3h+1$.\\
%If the first hop is like the one described in the second case.

%We again split $\mathbf{s}^{(1)}$ into two smaller strings $s^{(1)}_1,s^{(1)}_2$ which are the expansion of some string $s_1$ and $s_2$ respectively. We define $s^{(1)}_1$ as the first $p_i+2$ elements of $\mathbf{s}^{(1)}$, and $s^{(1)}_2$ be the elements of $\mathbf{s}^{(1)}$ starting from $p_i-1$ till the end. Again both $s^{(1)}_1$ and $s^{(1)}_2$ are both expansions of some strings as $p_i\equiv 2\pmod4$.\\
%Now the sequence $p_1,\dots p_i$ is a diffsequence in $s^{(1)}_1$ with $0$ hops, hence this corresponds to a diffsequence in $s_1$ with at most $1$ hop and of length at least $i-2$. Similarly, the sequence $p_{i+1},\dots p_{\ell}$ is a diffsequence in $s^{(1)}_2$ with $h-1$ hops; hence, this corresponds to a diffsequence in $s_2$ with at most $h$ hops and of length at least $\ell-i-3(h-1)-2$.\\ 
%Combining and subtracting one due to $s_{\text{pos}(p_i)}$ and $s_{\text{pos}(p_{i+1})}$ being the same (as this hop is actually not counted in $\mathbf{s}$). We get a diffsequence in $\mathbf{s}$ with at most $h+1$ hops of length at least $i-2+\ell-i-3(h-1)-2-1=\ell-3h-2$.
%\noindent

%We now explicitly construct a coloring of length $l\! \cdot \! 2^{2l}$ with the length of the maximum diffsequence at most $\frac{3l^2+3l}{2}$. Let $m_i$ denote the maximum length diffsequence with at most $i$ hops in $s_i$(recall that $s_i$ were the strings we constructed in Section~\ref{construction}). In particular, $m_l=l$. Now we can apply our lemma to bound $m_{i+1}-m_i$. Summing this, we can get a bound on $m_0$, which is the maximum length diffsequence in $s_0$. Now we claim \[m_i\le m_{i+1}+3i+2,\] for all $0\le i\le l-1$. To show this, let the maximum length diffsequence with at most $i$ hops in $s_i$ have $j\le i$ hops. Then by our lemma we have \[m_i\le x+3j+2\le m_{i+1}+3i+2,\]where $x$ is the maximum length diffsequence with at most $j+1\le i+1$ hops in $s_{i+1}$ which is at most $m_{i+1}$ due to maximality of $m_{i+1}$.
%\newline
%Therefore, \[m_0\le m_l+3\! \cdot \! \frac{l(l-1)}{2}+2l=\frac{3l^2+3l}{2}.\]



% We claim this is a $D$-diffsequence with at most $1$ hop in $\mathbf{s}$, and of length at least $\ell-2$, i.e., $\text{pos}(\mathcal{P}) \in \Psi_{\mathbf{s}}(1)$, and $|\text{pos}(\mathcal{P})| \ge l-2$.  If it exists, let $k$ be such that $s_{\text{pos}(p_i)}=0$ if and only if $i\le k$. Consider $p_1,\dots,p_k$, we know this is a diffsequence in $\mathbf{s}^{(1)}$ with $\textbf{s}^{(1)}_{p_i}=1$ and $s_{\text{pos}(p_i)}=0$ for all $i\le k$ therefore by Corollary~\ref{pos bound cor} we get the set $\{\text{pos}(p_i)\}_{1\le i\le k}$ has at least $k-1$ elements. Following a similar argument, we consider $p_{k+1},\dots,p_{\ell}$. Again this is a diffsequence in $\mathbf{s}^{(1)}$ and we also have $s^{(1)}_{p_i}=1$ and $s_{\text{pos}(p_i)}=1$ for all $k<i\le \ell$. Hence the set $\{\text{pos}(p_i)\}_{k< i\le \ell}$ has at least $\ell-k-1$ elements. Combining this with the above, we get that $\{\text{pos}(p_i)\}_{i=1}^{\ell}$ has at least $\ell-2$ distinct elements.
%Let $f:\{s_1,s_2,\! \cdot \!s,s_u\}\to\{A,B\}$ be a function such that, $\\
%f(s_i)=\begin{cases}A\quad \text{if the color at position } \lceil\frac{s_i}{4}\rceil \text{ in }T \text{ is 0}\\
%B\quad \text{if the color at position } \lceil\frac{s_i}{4}\rceil \text{ in }T \text{ is 1}\end{cases}\\$
%(basically $f$ says where does $s_i$ lie in the reduced string for $S$ that is $T$, our aim is to convert the diffsequence $\{s_i\}$ into a diffsequence in $T$, by squishing each $s_i$ into into the part it correspond in $T$, which formally is $\lceil\frac{s_i}{4}\rceil$)\\
%if $s_{i+1}-s_i\equiv0\pmod{4}$, then $\lceil\frac{s_{i+1}}{4}\rceil- \lceil\frac{s_i}{4}\rceil=\frac{s_{i+1}-s_i}{4}$ which is a power of $2\\$
%if $s_{i+1}-s_i=2$, then the only way this can happen is $001\color{red}1$ $\color{black}1\color{red}1\color{black}00$ or $00\color{red}1\color{black}1$ $\color{red}1\color{black}100$ in either case $\lceil\frac{s_{i}}{4}\rceil\neq\lceil\frac{s_{i+1}}{4}\rceil$ and $f(s_i)=A,f(s_{i+1})=B$ $\\$
%if $s_{i+1}-s_i=1$, then the only way this can happen is $00\color{red}11$ or $\color{red}11\color{black}00$ or $001\color{red}1$ $\color{red}1\color{black}100$. And only in the first and second case $\lceil\frac{s_{i}}{4}\rceil=\lceil\frac{s_{i+1}}{4}\rceil$, also note that $\lceil\frac{s_{i+1}}{4}\rceil\neq\lceil\frac{s_{i+2}}{4}\rceil$ as we cannot have such a structure that is followed by $i+1, i+2.$ Hence whenever $\lceil\frac{s_{i}}{4}\rceil=\lceil\frac{s_{j}}{4}\rceil$ we must have $j=i+1$ and $s_{i+1}-s_i=1,s_i\equiv2\text{ or } 0\pmod{4}$ $\\$
%Proposition 1 tells us that $\text{pos}(s_{i+1})-\text{pos}(s_i)$ is a power of whenever it is not 0. And as there cannot exist any $i$ for which $\text{find}(s_i)=B$ and $\text{find}(s_{i+1})=A$. Thus we assume $\text{find}(s_1),\text{find}(s_2),\dots,\text{find}(s_k)$ are all equal to $\mathbf{a}$ and the rest are $\mathbf{b}$ and $\text{pos}(s_{k+1})=\text{pos}(s_k)+1$. Therefore $X$ is a $D$-diffsequence with at most $1$ consecutive hop in $T$.
%We now show $|X|\ge u-2$. For this we need to remove the repeated elements in $\text{pos}(s_i)$, which is only possible if $\text{pos}(i)=\text{pos}(s_{i+1})$. Therefore we bound the number of times $\lceil\frac{s_{i}}{4}\rceil=\lceil\frac{s_{i+1}}{4}\rceil\\$
%Note that there is no index $i\leq k-1$ with $s_{i+1}-s_i=2$ (if there were then $\text{find}(s_{i+1})=B$ but $i+1\leq k$) .Therefore there is atmost one index $i\leq k-1$ with $\lceil\frac{s_{i}}{4}\rceil=\lceil\frac{s_{i+1}}{4}\rceil$ as otherwise if there were $2$ such indices, then let $i_1,i_2$ be the smallest such indices, then by the arguments in Proposition 1 we must have $s_{i_1}\equiv 2\pmod{4}$  as we assumed $i_1$ is the smallest index thus all  indices before $i_1$ must have the same residue modulo 4, and as the sequence starts with $\mathbf{a}$ we get $s_{i_1}\equiv2\pmod{4}$  and $s_{i_1+1}=s_{i_1}+1\equiv2+1\equiv3\pmod{4}$ and $s_{i_2}\equiv2\pmod{4}$ but by our assumption we must have $s_j\equiv s_{j+1}\pmod{4}$ for all $i_1+1\leq j\leq i_2$, a contradiction .$\\$
%Similarly there is atmost one index $k+1\leq i$ with $\lceil\frac{s_{i}}{4}\rceil=\lceil\frac{s_{i+1}}{4}\rceil$. Totally we get there are atmost two indices  $i$ with $\lceil\frac{s_{i}}{4}\rceil=\lceil\frac{s_{i+1}}{4}\rceil$, therefore $|X|\geq u-2\\$


%Let $X=\{t_1,t_2,\! \cdot \!s,t_v\}$ with $v\geq u-2$, and by the argument presented before, it is easy too see that $t_{i+1}-t_{i}$ is a power of $2$ (as it is never zero). Furthermore there is atmost one index $i$ with $t_{i+1}-t_i=1 $ and the colors at positions $t_i$ and $t_{i+1}$ in $T$ are different, Thus $X$ is a $D$-diffsequence with at most $1$ consecutive hop in $T$ with length $v\geq u-2$, as desired. $\\\\$

%Let $S_l$ be the string with $l$ digits and has alternating ones and zeroes, for all $l-1\geq i\geq 0$ define $S_i$ as the string obtained by replacing each $0$ in $S_{i+1}$ by $0011$ and $1$ by $1100$. Thus $S_i$ has length $4^l\! \cdot \! l$.
%. Let $m_i$ be length of the  maximum  diffsequence with at most $i$ hops in $S_i$. By our above theorem $m_i\leq m_{i+1}+3i+2$  thus $m_0\leq m_l+3\frac{l(l-1)}{2}+2l$ and note that $m_l=l$ (as $S_l$ has length $l$ ), therefore $m_0\leq \frac{3l^2+3l}{2}\\$


\begin{thebibliography}{9}
    \bibitem{Waerdan} 
        {\sc B. L. van der Waerdan.,}
        Beweis einer baudetschen vermutung. {\em Nieuw Archief voor Wiskude}, 15:212-216, 1927

    \bibitem{book}
        {\sc Landman B.} and {\sc Robertson A.,} Ramsey Theory on the integers. {\em Vol. 73. American Mathematical Soc.}, 2014.
    
    \bibitem{LR}
        {\sc Landman B.} and {\sc Robertson A.,} Avoiding monochromatic sequences with special gaps. {\em SIAM Journal on Discrete Mathematics}, 21.3 (2007): 794-801.

    \bibitem{Clifton}
        {\sc Clifton A.} New bounds on diffsequences. Discrete Mathematics. 2024 May 1;347(5):113929.
    \bibitem{Wesley}
        {\sc Wesley W. J.} Improved Ramsey-type theorems for Fibonacci numbers and other sequences,  arXiv preprint arXiv:2211.05167, 2022
    \bibitem{LV10}
        {\sc Landman B.} and {\sc Ventullo K.,}
        Avoiding monochromatic sequences with gaps in a fixed translation of the primes. Utilitas Mathematica, 82,207, 2010 207.
    \bibitem{CCLS18}
        {\sc Chokshi K.,} {\sc Clifton A.,}  {\sc Landman B.} and  {\sc Sawin O.}. "Ramsey functions for sequences with restricted gaps." Journal of Combinatorial Mathematics and Combinatorial Computing 105 (2018): 35-44.

\end{thebibliography}

\end{document}
